%!TEX root = ../template.tex
%%%%%%%%%%%%%%%%%%%%%%%%%%%%%%%%%%%%%%%%%%%%%%%%%%%%%%%%%%%%%%%%%%%%
%% abstract-pt.tex
%% NOVA thesis document file
%%
%% Abstract in Portuguese
%%%%%%%%%%%%%%%%%%%%%%%%%%%%%%%%%%%%%%%%%%%%%%%%%%%%%%%%%%%%%%%%%%%%

\typeout{NT FILE abstract-pt.tex}%


A Internet tornou-se uma ferramenta indispensável para aceder à informação e exercer a liberdade de expressão, um direito fundamental que dá poder aos indivíduos em todo o mundo. No entanto, os regimes autoritários e as agências governamentais tendem a suprimir este direito através de actividades de vigilância e censura. Estes esforços têm-se tornado cada vez mais sofisticados, empregando técnicas como ataques de análise de tráfego e métodos avançados de correlação de tráfego para monitorizar e minar o anonimato dos utilizadores.

Em resposta à crescente procura de comunicações seguras e privadas, as redes de anonimato, como o Tor, surgiram como ferramentas essenciais para salvaguardar as comunicações em linha de ponta a ponta e combater a censura, de que milhões de pessoas dependem diariamente. No entanto, investigações recentes demonstraram que potenciais ataques de desanonimização podem ser realizados por adversários poderosos a nível estatal através da correlação de tráfego e da recolha de impressões digitais utilizando técnicas avançadas de aprendizagem automática profunda.

Esta dissertação aborda estes desafios, reforçando as defesas da rede Tor contra os ataques acima referidos. O principal objetivo é desenvolver uma solução inovadora destinada a melhorar os nós de retransmissão Tor, tirando partido da arquitetura de software existente e introduzindo garantias formais de privacidade baseadas em circuitos privados diferenciados. Esperamos que a nossa solução proposta mantenha níveis práticos de débito e latência para a comunicação de ponta a ponta, garantindo a compatibilidade com cenários de utilização reais e proporcionando aos utilizadores uma proteção do anonimato adaptável e mais robusta.

Como já foi referido, o núcleo da solução proposta é a integração dos princípios da privacidade diferencial para introduzir dinamicamente ruído aleatório cuidadosamente delimitado nos fluxos de tráfego, atenuando os ataques de correlação e permitindo uma melhor solução resistente à recolha de impressões digitais. 
Tanto quanto sabemos, os contributos esperados para a dissertação podem representar um esforço pioneiro na incorporação da Privacidade Diferencial na arquitetura de software dos retransmissores Tor.
Diferencial na arquitetura de software dos nós de retransmissão Tor. Além disso, a utilização da Privacidade Diferencial permitirá a criação de um mecanismo formalmente comprovado, escalável e eficaz que reforça o anonimato online, resistente às técnicas mais avançadas de censura na Internet e que defende o direito fundamental à liberdade de expressão num mundo digital cada vez mais monitorizado.


\keywords{
    Redes de Anonimato
    \and Privacidade Diferencial
    \and Comunicação Preservadora de Privacidade
    \and Ataques de Análise de Tráfego
    \and Garantias Reforçadas de Privacidade Ponto a Ponto 
    \and Correlação de Tráfego
}

